% Wrapper-free proof; uses only amsmath, amssymb and theorem environments.
\section{Four retained coordinate candidates and the exact residual-three first hit}
\label{sec:boundary-four-candidates}

The antecedent \emph{Signed margins, exact witness swaps, and the
prime-output sieve}, Theorem \texttt{thm:margin-swap-bands}, proves
the same-tag exchange, its shell criterion, its oriented divisor map,
and the initial forbidden bands. The source file is
\texttt{signed\_margin\_sieve.tex}, SHA256
\texttt{d51347ad95517c59d4922b526bc55ac1e7a9487f18503f415943b9870ddda6ef}.
We compute here the two additional tag candidates, every failed
candidate's exact rational obstruction, the complete integral orbit
fibres, and the actual least code on every occupied residual-three
shell. All raw coordinates and their source marks remain present.

Fix a prime \(p=12h+1\), \(h\ge1\), and a full original shell
\[
 0\le j<6h,\qquad a=3h+j+1,\qquad R=4j+3=4a-p,
 \qquad S=pa.
\]
Let \(\mathcal R_{p,j}\) consist of the complete records
\[
\begin{gathered}
 X=(p,j,a,R,S,A,B,C,D,\varepsilon),\\
 A,B,D\in\mathbb Z_{>0},\quad ABD=a,\\
 \varepsilon\in\{0,1\},\quad
 C=\frac{A+p^{1-\varepsilon}B}{R}\in\mathbb Q_{>0}.
\end{gathered}
\]
Let \(\mathcal I_{p,j}\) be its exact subset \(C\in\mathbb Z\).
No coprimality restriction is imposed on \(A,B\) in these raw
domains. Since \(a<p\), \(\gcd(a,R)=\gcd(a,p)=1\);
therefore every one of \(A,B,D,p\) is a unit modulo \(R\).
For the first equality, a prime dividing \(a,R\) would divide
\(p=4a-R\), contradicting \(a<p\). For the second use primality
and \(0<a<p\). Finally \(R\) is odd and
\(\gcd(p,R)=1\), since \(p\mid R\) would imply \(p\mid a\).
These unit statements will justify each cancellation explicitly.

\begin{theorem}[The four rational candidates with their complete inverse]
\label{thm:boundary-four-rational}
For \(s,t\in\{0,1\}\), set
\[
 (\alpha,\beta)=
 \begin{cases}(A,B)&s=0,\\(B,A)&s=1,\end{cases}
 \qquad \eta=\varepsilon+t-2\varepsilon t,
 \qquad C_{s,t}=\frac{\alpha+p^{1-\eta}\beta}{R}.
\]
The map \(T_{s,t}:\mathcal R_{p,j}\to\mathcal R_{p,j}\)
retains \(p,j,a,R,S,D\) and replaces
\((A,B,C,\varepsilon)\) by
\((\alpha,\beta,C_{s,t},\eta)\). Its inverse is precisely
\(T_{s,t}\). The four maps form an action of
\((\mathbb Z/2\mathbb Z)^2\): composition adds both displayed
bits modulo two. Every individual map has singleton fibres.
The full rational orbit of \(X\) has four distinct records if
\(A\ne B\), and two if \(A=B\); its complete stabilizer is
\(\{(0,0)\}\) in the former case and
\(\{(0,0),(1,0)\}\) in the latter case.

For every candidate retain its original ordered rational witness and
both oriented residuals:
\begin{align*}
 (x_{s,t},y_{s,t},z_{s,t})
   &=(a,p^\eta\alpha C_{s,t}D,p\beta C_{s,t}D),\\
 \frac1{x_{s,t}}+\frac1{y_{s,t}}+\frac1{z_{s,t}}&=\frac4p,\\
 d_{y,s,t}=Ry_{s,t}-S&=p^\eta\alpha^2D,\\
 d_{z,s,t}=Rz_{s,t}-S&=p^{2-\eta}\beta^2D,
 \qquad d_{y,s,t}d_{z,s,t}=S^2.
\end{align*}
Thus these residuals are positive integers even when one or both
witness denominators are nonintegral. Writing
\(u=B^2D\), the retained target coordinate is
\(u_{s,t}=u\) if \(s=0\), and
\(u_{s,t}=a^2/u=A^2D\) if \(s=1\), with the tag
\(\eta\) still attached. The map to the graph containing the
source \(X\), both bits, and this full target has inverse projection
to \((X,s,t)\), hence singleton fibres; no witness ordering or
choice among equal rational candidates is forgotten.

Set \(k=\gcd(A,B)\). The raw-coordinate correspondence is
\[
 (A_0,B_0,C_0,D_0,k)
   =(A/k,B/k,C/k,k^2D,k),\qquad
 (A,B,C,D)=(kA_0,kB_0,kC_0,D_0/k^2).
\]
Here \(\gcd(A_0,B_0)=1\), \(k^2\mid D_0\), and
\(C_0=(A_0+p^{1-\varepsilon}B_0)/R\).
The maps \(T_{s,t}\) commute with this correspondence, retaining
the very same \(k\) and \(D_0\). Above a fixed coprime record,
the full raw-scale fibre is exactly
\(\{k>0:k^2\mid D_0\}\), including when \(C_0\) is rational.
The integral locus is equivalent to \(C_0\in\mathbb Z\), and
on that locus all displayed witnesses agree under this raw-scale
correspondence coordinate by coordinate.
\end{theorem}
\begin{proof}
The product \(\alpha\beta D=ABD=a\) and the target coefficient
equation put each target in the specified rational domain. Applying
the same bit twice restores \(A,B,\varepsilon\). The coefficient
is the unique rational number satisfying its defining equation, so
it also returns to the original \(C\). The swap and tag operations
commute on these three discrete coordinates, which proves the group
law and all individual inverses. A fixed point must have \(t=0\),
because changing the tag changes that retained coordinate. For
\(s=1,t=0\) it must, and conversely does, have \(A=B\).
These are all four possible stabilizer elements, giving the orbit
counts directly.

The target equation reads
\[
 4\alpha\beta C_{s,t}D
     =\alpha+p^{1-\eta}\beta+pC_{s,t}.
\]
Over the positive rational common denominator
\(p\alpha\beta C_{s,t}D\), the three reciprocals have
numerators \(pC_{s,t},p^{1-\eta}\beta,\alpha\), respectively.
Their sum is consequently \(4/p\). Multiplying
\(RC_{s,t}=\alpha+p^{1-\eta}\beta\) by
\(p^\eta\alpha D\), and then by \(p\beta D\), gives
the two residual identities after subtracting
\(S=p\alpha\beta D\). Their product is
\(p^2\alpha^2\beta^2D^2=S^2\). The formulas for
\(u_{s,t}\) follow from \(a^2=A^2B^2D^2\); applying the
same complement twice restores \(u\), still with its appropriate
tag. Projection from the fully specified graph restores exactly
the source and the two chosen bits, including at fixed points.

The raw-scale formulas are inverse because
\(\gcd(kA_0,kB_0)=k\). Substitution gives
\(A_0B_0D_0=a\) and the stated rational coefficient. The
square-divisor condition is exactly what makes \(D_0/k^2\)
an integer. Thus every scale in the displayed fibre occurs, and
no other scale does. On the original integral locus, reducing
\(RC=A+p^{1-\varepsilon}B\) modulo \(k\) gives
\(k\mid RC\); since \(k\mid a\), \(\gcd(k,R)=1\), so
\(k\mid C\). This proves \(C_0\) integral. Conversely,
\(C=kC_0\) is integral when \(C_0\) is. The same argument
works for every target because its \(k\) is unchanged.
Finally direct substitution gives
\(ABD=A_0B_0D_0\),
\(p^\varepsilon ACD=p^\varepsilon A_0C_0D_0\), and
\(pBCD=pB_0C_0D_0\). The target calculation has the same
identities with \(\alpha/k,\beta/k,C_{s,t}/k\). This proves
the asserted commutation and every retained witness equality.
\end{proof}

\begin{theorem}[Exact obstruction classes, denominators, and four integral fibres]
\label{thm:boundary-obstruction-fibres}
For an integral source \(X\in\mathcal I_{p,j}\), the following
table gives the target numerator modulo \(R\), the common exact
positive reduced denominator \(\delta_{s,t}\) of
\(C_{s,t},y_{s,t},z_{s,t}\), and the complete integrality test.
The original positive rational numbers themselves remain the
fractions defined in the preceding theorem.
\[
\begin{array}{c|c|c|c}
 (s,t)\text{ and source tag}&
 \alpha+p^{1-\eta}\beta\pmod R&\delta_{s,t}&
 T_{s,t}X\in\mathcal I_{p,j}\\\hline
 (0,0),\ \varepsilon=0\text{ or }1&0&1&\text{always}\\
 (1,0),\ \varepsilon=1&0&1&\text{always}\\
 (1,0),\ \varepsilon=0&(1-p^2)B&
 R/\gcd(R,p^2-1)&R\mid p^2-1\\
 (0,1)\text{ or }(1,1),\ \varepsilon=0&(1-p)B&
 R/\gcd(R,p-1)&R\mid p-1\\
 (0,1),\ \varepsilon=1&(p-1)B&
 R/\gcd(R,p-1)&R\mid p-1\\
 (1,1),\ \varepsilon=1&-(p-1)B&
 R/\gcd(R,p-1)&R\mid p-1
\end{array}
\]
Each integral restriction
\[
 T_{s,t}:\{X\in\mathcal I_{p,j}:T_{s,t}X\in\mathcal I_{p,j}\}
   \longrightarrow
 \{X\in\mathcal I_{p,j}:T_{s,t}X\in\mathcal I_{p,j}\}
\]
is a bijection with the same inverse \(T_{s,t}\) and singleton
fibres. If the test fails, the precise rational target and its
inverse in \(\mathcal R_{p,j}\) still apply.

More explicitly, the additive morphism
\[
 \psi_R:\mathbb Z/R\mathbb Z\longrightarrow
      (\mathbb Q/\mathbb Z)[R],\qquad
 \overline n\longmapsto n/R+\mathbb Z
\]
is an isomorphism. The target obstruction is exactly
\(\omega_{s,t}=\psi_R(\overline{\alpha+p^{1-\eta}\beta})
 =C_{s,t}+\mathbb Z\), and has additive order
\(\delta_{s,t}\). For any integer multiplier \(m\) occurring
in the second column, put \(g_m=\gcd(R,m)\) and
\(\delta_m=R/g_m\), with \(\gcd(R,0)=R\).
The map \(\psi_R\circ[m]\) has image exactly
\((\mathbb Q/\mathbb Z)[\delta_m]\), kernel exactly
\(\delta_m\mathbb Z/R\mathbb Z\), and every nonempty fibre
is exactly \(\overline{b_0}+\delta_m\mathbb Z/R\mathbb Z\).
In the table the retained input \(B\) is a unit modulo \(R\),
so its output has the full order \(\delta_m\).
\end{theorem}
\begin{proof}
For \(\varepsilon=0\), the integral source gives
\(A\equiv-pB\pmod R\). Its same-tag swap numerator is
\(B+pA\equiv(1-p^2)B\). Both tag-changing numerators are
\(A+B\equiv(1-p)B\). For \(\varepsilon=1\), the source
gives \(A\equiv-B\pmod R\); its same-tag swap has numerator
\(A+B\equiv0\), while its tag changes have numerators
\(A+pB\equiv(p-1)B\) and
\(B+pA\equiv-(p-1)B\). The identity target has its original
integral numerator. This proves every residue entry independently
of any representative choice.

For an integer \(N>0\), the exact denominator of \(N/R\)
is \(R/\gcd(R,N)\). To prove this statement, write
\(N=gN_1,R=gR_1\), where \(g=\gcd(R,N)\) and
\(\gcd(N_1,R_1)=1\); then \(dN/R\) is an integer exactly
when \(R_1\mid d\), by Bezout. Its least positive such \(d\)
is \(R_1\). Every factor \(B,p,\alpha,\beta,D\) is a
unit modulo \(R\), so
\[
 \gcd(R,mB)=\gcd(R,m),\quad
 \gcd(R,p^\eta\alpha DN)=\gcd(R,N),\quad
 \gcd(R,p\beta DN)=\gcd(R,N).
\]
Apply the denominator calculation to each displayed numerator and
its two indicated multiples. This proves the same exact denominator
for the coefficient and both ordered witness denominators, all table
values, and both directions of every integrality assertion.
The restricted map is its own inverse because the rational map is
its own inverse; if its source and target are integral, the same
condition holds with their roles exchanged.

Changing \(n\) by a multiple of \(R\) changes \(n/R\) by
an integer, so \(\psi_R\) is well-defined and additive. Its
kernel is zero. Every class annihilated by \(R\) is represented
by a rational \(x\) with \(Rx=n\in\mathbb Z\), and is
therefore \(\psi_R(\overline n)\). This proves its stated
codomain and both parts of the isomorphism. The denominator
calculation already proves the order of each \(N/R+\mathbb Z\).
For \(m\ne0\), write \(R=g_m\delta_m\),
\(m=g_mm_1\), \(\gcd(\delta_m,m_1)=1\).
Then \(R\mid mb\) is equivalent to \(\delta_m\mid b\),
giving the complete kernel. Bezout gives an integer inverse of
\(m_1\) modulo \(\delta_m\), so the fractions
\(mb/R=m_1b/\delta_m\) attain every class
\(n/\delta_m+\mathbb Z\); these are precisely the stated
image. Two inputs have the same output exactly when their
difference is in the proved kernel, giving every fibre. For
\(m=0\), the map is zero, \(\delta_m=1\), and these same
descriptions respectively give the zero image, whole domain, and
whole nonempty fibre. Finally \(\gcd(B,R)=1\) gives
\(\gcd(B,\delta_m)=1\), so multiplication by the retained
\(B\) leaves the exact order \(\delta_m\) unchanged.
\end{proof}

\begin{theorem}[Complete integral orbits and the original code-lowering map]
\label{thm:boundary-orbit-code}
Retain either original radix
\[
 (J_H,M_H)=(3h,6h),\qquad (J_F,M_F)=(6h,9h),\qquad
 \Lambda_\rho=2J_\rho M_\rho.
\]
For \(\rho=H\), impose its actual shell domain \(j<3h\);
for \(\rho=F\), keep the whole domain above. Codes use the
retained coprime coordinates of the raw record:
\[
 c_\rho(X)=(A_0-1)\Lambda_\rho+2M_\rho j
                    +2(B_0-1)+\varepsilon.
\]
The rectangle \(1\le A_0,B_0\le M_\rho\),
\(0\le j<J_\rho\), \(\varepsilon\in\{0,1\}\) has
the following full code inverse on
\(0\le c<M_\rho\Lambda_\rho\):
\[
\begin{gathered}
 A_0=1+\lfloor c/\Lambda_\rho\rfloor,\quad
 r=c-(A_0-1)\Lambda_\rho,\quad
 j=\lfloor r/(2M_\rho)\rfloor,\\
 w=r-2M_\rho j,\quad B_0=1+\lfloor w/2\rfloor,\quad
 \varepsilon=w-2\lfloor w/2\rfloor.
\end{gathered}
\]
On a hit this recovers \(a=3h+j+1\), \(R=4j+3\),
\(S=pa\), \(D_0=a/(A_0B_0)\), and
\(C_0=(A_0+p^{1-\varepsilon}B_0)/R\).
Keeping the scale \(k\) then recovers the entire raw record.
Every integral target in the same rational orbit has code
\begin{equation}\label{eq:boundary-four-code-difference}
 c_\rho(T_{s,t}X)-c_\rho(X)
  =s(B_0-A_0)(\Lambda_\rho-2)+(\eta-\varepsilon).
\end{equation}
Let \(\mathcal O_I(X)\) be the rational orbit intersected with
the integral domain. Its complete list and its smallest-code member
are as follows; the coefficient at every listed member is always
its forced coefficient and \(D,k,p,j,a,R,S\) are retained.
\[
\begin{array}{c|c|c}
\text{shell and tag}&\text{members, as }(A,B,\varepsilon)&
                 \text{smallest-code coordinates}\\\hline
 R\mid p-1& (A,B,0),(B,A,0),(A,B,1),(B,A,1)&
               (\min(A,B),\max(A,B),0)\\
 R\nmid p-1,\ R\mid p^2-1&
              (A,B,\varepsilon),(B,A,\varepsilon)&
               (\min(A,B),\max(A,B),\varepsilon)\\
 R\nmid p^2-1,\ \varepsilon=1&(A,B,1),(B,A,1)&
               (\min(A,B),\max(A,B),1)\\
 R\nmid p^2-1,\ \varepsilon=0&(A,B,0)&(A,B,0)
\end{array}
\]
The first row has four distinct members. The second row has two
except when \(A=B\), when it has one. The third row has two
distinct members. The fourth row has one. The map selecting the
listed smallest member, while retaining \(k\), has full fibre
exactly that displayed integral orbit. If \(k\) is also forgotten,
the full fibre above its coprime smallest member is the disjoint
union of these orbits over every \(k>0\) with \(k^2\mid D_0\).
Retaining the source and chosen bits instead gives the singleton
graph inverse of Theorem~\ref{thm:boundary-four-rational}.

Consequently any globally least hit in either original radix lies
among the listed orbit minima. In particular on a shell
\(R\mid p-1\) its tag is zero and \(A_0<B_0\).
On every middle hit the swap lowers the original code exactly when
\(A_0>B_0\). On an exterior hit with \(A_0>B_0\), the
same-tag swap gives a lower integral hit exactly on
\(R\mid p^2-1\); elsewhere
the rational companion and its nontrivial obstruction remain the
exact relation. These statements specify both the retained positive
integer domains and the failed integral fibres.
\end{theorem}
\begin{proof}
The product \(A_0B_0D_0=a\) gives
\(1\le A_0,B_0\le a\le M_\rho\) on the specified shell
range. Swapping retains these same bounds. Thus all candidates
remain in the same code rectangle. The inner digit
\(2(B_0-1)+\varepsilon\) runs bijectively through
\(0,\ldots,2M_\rho-1\). Adding \(2M_\rho j\) runs
bijectively through \(0,\ldots,\Lambda_\rho-1\).
Adding \((A_0-1)\Lambda_\rho\) therefore runs bijectively
through the full displayed interval. Successive Euclidean division
gives exactly the three remainders \(r,w,\varepsilon\) above,
so both encode--decode composites are identities on every original
digit. The hit's defining budget and gate make the displayed
\(D_0,C_0\) positive integers; their formulas and the raw-scale
inverse restore every other recorded coordinate.
The tag increment is precisely
\(\eta-\varepsilon\); the difference of the outer and inner
digits is \((B_0-A_0)\Lambda_\rho+2(A_0-B_0)\) if
\(s=1\), and zero otherwise. This proves the full difference
formula without changing the radix or any shell digit.

The four-candidate table of the preceding theorem gives exactly the
four rows of integral members. If \(R\mid p-1\), equality
\(A=B\) would imply, from either source tag, \(R\mid2A\).
But \(A\) and two are units modulo the odd integer \(R\ge3\),
a contradiction. A middle hit with \(A=B\) gives that same
contradiction even without \(R\mid p-1\). In the second row
an exterior hit with \(A=B\) is permitted, and the swap then
fixes the entire record. There are no other coincidences, by the
stabilizer calculation above.

One has \(\Lambda_\rho\ge36\), so
\(\Lambda_\rho-2>1\). Thus any nonzero integer decrease of
\(A_0\) under a swap decreases the code by more than either
possible tag increment; at a fixed ordered pair tag zero precedes
tag one by exactly one. This proves each claimed smallest member
and its uniqueness as a record. Distinct integral orbit minima
cannot share a fibre: if two records have the same minimum, they
belong to the rational orbit of that minimum. Conversely every
integral member of that orbit has precisely the same computed
minimum. This proves the complete fibre. The raw-scale inverse
already established gives exactly the extra disjoint choices when
\(k\) is forgotten, and proves that no other raw preimages occur.

If a globally least hit were not minimal within its own integral
orbit, the computed smaller member would be another original hit
with a strictly smaller original code, a contradiction. All further
statements now follow from the displayed minima, the exclusion of
\(A=B\) in the first row, and the exact difference formula.
\end{proof}

\begin{theorem}[The exact global first code on the occupied residual-three domain]
\label{thm:boundary-r3-global-minimum}
Put \(a_3=3h+1=(p+3)/4\), and retain its actual prime factors.
Define the explicit prime set
\[
 \mathcal P_2(a_3)=\{\ell:\ell\text{ prime},\ \ell\mid a_3,
                                  \ \ell\equiv2\pmod3\}.
\]
The residual-three shell has an integral hit exactly on the domain
\(\mathcal P_2(a_3)\ne\varnothing\). On this precisely
specified domain let \(\ell_*=\min\mathcal P_2(a_3)\).
The global least hit in each of the two original radices is
\[
 (A_0,j,B_0,\varepsilon)=(1,0,\ell_*,0),\qquad
 c_{*,H}=c_{*,F}=2(\ell_*-1),
\]
with every original coordinate given by
\[
 D_0=a_3/\ell_*,\quad C_0=(1+p\ell_*)/3,\quad
 u=a_3\ell_*,\quad R=3,\quad S=pa_3,
\]
\[
 (x,y,z)=(a_3,C_0D_0,p\ell_*C_0D_0),\qquad
 d_y=D_0,\quad d_z=p^2\ell_*^2D_0.
\]
Its complete raw-scale fibre is
\[
 (A,B,C,D,\varepsilon)
    =(k,k\ell_*,kC_0,D_0/k^2,0),\qquad k>0,\quad k^2\mid D_0.
\]
If \(\mathcal P_2(a_3)=\varnothing\), the entire original
\(j=0\) hit fibre is empty in both radices; the theorem makes
no assertion that the later shells are empty.

For a passing boundary source with coprime digits
\((A_0,j,B_0,\varepsilon)=(m+1,0,1,0)\), its same-tag
swap has code \(2m\), retaining
\[
 D_0=a_3/(m+1),\quad
 C_0=(m+1+p)/3,\quad C_0^{\leftrightarrow}=(1+p(m+1))/3.
\]
The global formula strengthens this to
\(c_{*,\rho}=2(\ell_*-1)\le2m\), with equality exactly
when \(m+1=\ell_*\). The swapped ordered witness is
\((a_3,C_0^{\leftrightarrow}D_0,
 p(m+1)C_0^{\leftrightarrow}D_0)\), and its source is
recovered by the proved involution with the same raw scale.
\end{theorem}
\begin{proof}
Here \(a_3\equiv1\pmod3\), so none of its prime factors is
three. For any raw hit, its coprime coordinates satisfy
\(A_0B_0\mid a_3\) and, because \(p\equiv1\pmod3\),
the exact gate \(A_0+B_0\equiv0\pmod3\), independently
of the retained tag. Both coordinates are units modulo three;
therefore at least one of \(A_0,B_0\) is two modulo three.
Its prime factorization contains a prime two modulo three: if
every factor were one, its product would be one. That prime divides
\(a_3\), proving necessity and proving that no omitted raw scale
can create a hit on the failed domain.

Conversely any \(\ell\in\mathcal P_2(a_3)\) gives
\(A_0=1,B_0=\ell,D_0=a_3/\ell\) and
\(C_0=(1+p\ell)/3\), all positive integers. Their equation
\(3C_0=1+p\ell\), together with \(4a_3=p+3\), is exactly
the original central equation. The preceding witness calculation
therefore proves it is a hit with the displayed ordering, residuals,
and full scale fibre. This establishes the occupied-domain
equivalence before making any least-code assertion.

On this domain select \(\ell_*\). Every hit with \(A_0\ge2\)
has code at least \(\Lambda_\rho\). A hit with \(A_0=1\)
and \(j\ge1\) has code at least \(2M_\rho\). The selected
hit has code
\(2(\ell_*-1)\le2(a_3-1)<2M_\rho\le\Lambda_\rho\),
where \(a_3\le M_\rho\) and \(J_\rho\ge1\).
It therefore precedes all these other digits. The remaining hits
have \(A_0=1,j=0\), hence \(B_0\mid a_3\) and
\(B_0\equiv2\pmod3\). Such a positive integer has at least
one prime factor in \(\mathcal P_2(a_3)\), so
\(B_0\ge\ell_*\). Equality at the smallest code is attained
with tag zero, and tag one is one integer later. Thus no code
precedes \(2(\ell_*-1)\), and the decoder makes its digits
unique. This proves the actual global minimum in both rectangles,
including comparison with every later shell.

For the boundary source, its budget says \(m+1\mid a_3\)
and its gate says \(m+1\equiv2\pmod3\). In particular
\(m+1\ge2\), and it has a prime factor at least \(\ell_*\)
and at most \(m+1\). The shell has \(3\mid p-1\), so its
integral swap and both coefficient formulas follow from the
four-candidate computation. Substitution into the unchanged radix
gives the exact swapped code \(2m\). The global formula then
gives the inequality, and equality of the two displayed integer
codes is equivalent to \(\ell_*=m+1\). The witness and inverse
are the specified candidate formulas and preserve every raw scale.
\end{proof}

\paragraph{Two exact residual-seven companions.}
For \(R=7\), the exterior same-tag integral restriction is
precisely \(p\equiv1\) or \(6\pmod7\): the prime seven
divides \((p-1)(p+1)\) exactly when it divides one of the two
factors. The tag-changing restriction is precisely
\(p\equiv1\pmod7\). The maps and fibres in the preceding
theorems apply on each of these exact residue domains.
At \(p=373\), \(a=95\), \(R=7\), the two original exterior
records \((A,B,C,D,\varepsilon)=(19,1,56,5,0)\) and
\((5,1,54,19,0)\) have \(p\equiv2\pmod7\).
Their same-tag swaps have respectively
\[
 C^{\leftrightarrow}=7088/7,\quad
 (x',y',z')=(95,35440/7,251163280/7),
\]
and
\[
 C^{\leftrightarrow}=1866/7,\quad
 (x',y',z')=(95,35454/7,66121710/7).
\]
These follow by substituting the two retained source records in
\(C^{\leftrightarrow}=(B+pA)/R\),
\(y'=BC^{\leftrightarrow}D\), and
\(z'=pAC^{\leftrightarrow}D\). Since
\(\gcd(7,p^2-1)=\gcd(7,3)=1\), the theorem proves that
every displayed coefficient and denominator has exact denominator
seven and that both rational witnesses have reciprocal sum \(4/373\).
The tag-changing candidates also have denominator seven because
\(\gcd(7,p-1)=1\). Thus the complete integral orbit of each
of these exterior records is its singleton, while all three other
labelled candidates and their inverse maps remain fully specified.
For completeness, trial division proves the asserted primality of
373: the remainders at \(2,3,5,7,11,13,17,19\) are respectively
\(1,1,3,2,10,9,16,12\), all nonzero; these are all primes below
\(\sqrt{373}<20\).
